% 2de graad arbeidsmarktfinaliteit, Wiskunde in de praktijk (rekenen, 2D en 3D, grafieken en tabellen lezen).
% Monet: een werkhaven bij zonsopgang, met kranen, schepen en schoorstenen.
% Motief: een staafdiagram met het gemiddelde, de stippellijn waar de staven samen op uitkomen.
\def\omslagbestand{omslag-impressie.jpg}
\def\omslagbron{Claude Monet, \textit{Impressie, zonsopgang}, 1872 (publiek domein)}
\def\omslagrand{10.9 + 0.22*sin(\x*360/6.5)}
\def\omslagmotief{%
  \randlijn
  \foreach \plek/\hoogte in {13.6/2.2,15.2/3.6,16.8/2.8,18.4/4.6}{%
    \fill[white,fill opacity=.17] (\plek,4.2) rectangle ++(1.1,\hoogte);
    \draw[white,line width=.9pt] (\plek,4.2) rectangle ++(1.1,\hoogte);
  }
  \draw[white,line width=.8pt] (13.2,4.2) -- (19.9,4.2);
  \draw[fel,line width=1.2pt,dashed] (13.2,7.5) -- (19.9,7.5);
  \node[fill=fel,text=inkt,rounded corners=1.5pt,inner sep=2.5pt,font=\kopletter\small,anchor=south west] at (13.2,8.95) {gemiddelde};
  \draw[fel,line width=.8pt] (13.9,8.95) -- (13.9,7.5);
}
